% !TEX root =  mosaic.tex
\section{Overview of Mosaic}
\label{sec:overview}

\name treats the transformed program as the boundary between transformation
scheduling and execution-space optimization: existing schedulers fix the legal
computation order, and \name decides how that order should consume GPU
architectural resources, emitting tile sizes \(\boldsymbol{\tau}\),
thread-block dimensions \(\mathbf{b}\), and a shared memory placement scheme
\(S\). Throughout, \name operates outside the compiler; the backend is driven
as a black box and all resource reasoning happens in \name. As summarized in
\autoref{fig:mosaic-pipeline}, the workflow comprises four components:

\begin{itemize}
  \item Transformed-program acquisition. \name starts from a program
  whose schedule is already determined by an external compiler. Its loop
  transformations, including fusion, interchange, skewing, distribution, and
  related restructurings, are reflected in the schedule metadata and are neither
  re-derived nor overridden; \name reads schedule-tree annotations for PPCG and
  schedule traces for TVM.
  \item Metadata extraction and target profiling. \name extracts metadata for
  loops, accesses, operations, coalescing, and placement candidates, then
  combines it with once-per-architecture target limits and effective costs,
  yielding the finite execution-space problem instantiated in \autoref{sec:model}.
  \item Analytical execution-space solving. \name enumerates
  shared memory placement schemes for read-only operands and, for each scheme,
  solves over tile sizes and thread-block dimensions; feasibility is checked by
  hardware limits and the conservative register filter, and ranking uses the
  quantized wave cost with an LP tie-break among equal wave costs.
  \item Backend emission. The selected tile sizes, block sizes, and
  placement choices are emitted as backend parameters or CUDA macros.
\end{itemize}

We cast the workflow as a joint minimization over the execution-space candidate
set:
\begin{equation}
(\boldsymbol{\tau}^{*}, \mathbf{b}^{*}, S^{*})
 = \arg\min_{\substack{S \in \Sigma \\
(\boldsymbol{\tau}, \mathbf{b}) \in \mathcal{V}(S)}}
  T_{\mathit{total}}(\boldsymbol{\tau}, \mathbf{b}; S).
\label{eq:two-level}
\end{equation}
Here \(S\) ranges over the power set
\(\Sigma = 2^{\Sigma_{\mathit{cand}}}\) of read-only placement candidates and
\((\boldsymbol{\tau}, \mathbf{b})\) ranges over the placement-dependent feasible
domain \(\mathcal{V}(S)\). \autoref{eq:two-level} is a \emph{single} joint
minimization over the Cartesian product of placements and tile/block choices:
\name solves the inner problem to obtain the best
\((\boldsymbol{\tau}, \mathbf{b})\) cost for \emph{every} placement
\(S \in \Sigma\) and returns the lowest-cost triple overall.
\autoref{sec:metadata} describes the metadata and target profile that
instantiate \autoref{eq:two-level}; \autoref{sec:model} formalizes the cost
model, constraints, and solver.

For the running GEMM example, \name extracts
\(\mathcal{P} = \{i, j\}\), \(\mathcal{R} = \{k\}\), read-only operands \(A\),
\(B\), and register accumulator \(C\), enumerates the induced placement schemes,
solves for tile and block choices, and emits the selected parameters as CUDA
macros.


\begin{figure*}[tbp]
\centering
\resizebox{\textwidth}{!}{%
\begin{tikzpicture}[
  font=\scriptsize,
  >={Stealth[length=2mm]},
  box/.style   ={rectangle, rounded corners=2pt, draw=black!70,
                 align=left, inner sep=2.8pt, minimum height=7mm},
  input/.style ={box, fill=blue!5, text width=34mm},
  target/.style={box, fill=green!6, text width=34mm},
  stage/.style ={box, fill=orange!8, align=center, text width=20mm,
                 minimum height=9mm},
  output/.style={box, fill=purple!6, text width=31mm},
  title/.style ={font=\scriptsize\bfseries, align=center},
  arr/.style   ={->, thick, draw=black!75},
  darr/.style  ={->, thick, dashed, draw=black!55},
  dline/.style ={thick, dashed, draw=black!55},
]

\node[input] (sched) {%
\textbf{Schedule metadata}\\
\textit{Loop structure}\\
\quad loop variables and extents\\
\quad schedule bands\\
\quad tileable axes\\[1pt]
\textit{Data layout}\\
\quad array roles\\
\quad access layouts};

\node[target, below=3mm of sched] (prof) {%
\textbf{Target profile}\\
\textit{Hardware limits}\\
\quad registers, shared memory\\
\quad blocks, threads, SM count\\[1pt]
\textit{Effective costs}\\
\quad global/shared/register\\
\quad access costs\\
\quad synchronization and compute};

\begin{scope}[on background layer]
  \node[fit=(sched)(prof), inner sep=0pt] (inputgroup) {};
\end{scope}

\node[title, anchor=south]
      (inputtitle)
      at ([yshift=1.2mm]inputgroup.north)
      {External compiler / target inputs};

\node[stage, anchor=west]
      (cls)
      at ([xshift=7mm,yshift=7mm]inputgroup.east)
      {Dimension\\classification};

\node[stage, right=3mm of cls]
      (plc)
      {Placement\\schemes \(\Sigma\)};

\node[stage, right=3mm of plc]
      (cand)
      {Tile/block\\candidate sets};

\node[stage, below=5mm of cls]
      (feas)
      {Feasible domain\\\(\mathcal{V}(S)\)};

\node[stage, right=3mm of feas]
      (cost)
      {Cost model\\\(\phi,T_{\mathit{block}}\)};

\node[stage, right=3mm of cost]
      (solve)
      {MILP + LP\\selection};

\begin{scope}[on background layer]
  \node[
    draw=black!55,
    rounded corners=3pt,
    fill=orange!3,
    inner xsep=4mm,
    inner ysep=4mm,
    fit=(cls)(plc)(cand)(feas)(cost)(solve)
  ] (engine) {};
\end{scope}

\node[title, anchor=south]
      at ([yshift=1.2mm]engine.north)
      {\textsc{Mosaic} analytical engine
       \quad (for each \(S\in\Sigma\))};

\draw[arr] (cls.east) -- (plc.west);
\draw[arr] (plc.east) -- (cand.west);

\draw[arr, rounded corners=1.5pt]
  (cand.south) -- ++(0,-2mm) -| (feas.north);

\draw[arr] (feas.east) -- (cost.west);
\draw[arr] (cost.east) -- (solve.west);

\coordinate (schedexit) at (sched.east |- cls.west);
\draw[arr] (schedexit) -- (cls.west);

\coordinate (profileleft)
  at ([yshift=-2.5mm]engine.south west);

\coordinate (profilefeas)
  at (feas.south |- profileleft);

\coordinate (profilecost)
  at (cost.south |- profileleft);

\coordinate (profexit)
  at (prof.east |- profileleft);

\draw[dline] (profexit) -- (profilecost);
\draw[darr] (profilefeas) -- (feas.south);
\draw[darr] (profilecost) -- (cost.south);

\matrix (backendboxes) [
  matrix,
  anchor=west,
  row sep=3mm,
  inner sep=0pt
] at ([xshift=7mm]engine.east |- inputgroup.center) {

  \node[output] (out) {%
    \textbf{Selected execution space}\\
    placement \(S^{*}\)\\
    tile sizes \(\boldsymbol{\tau}^{*}\)\\
    block sizes \(\mathbf{b}^{*}\)\\
    register factors \(\mathbf{R}^{*}\)}; \\

  \node[output] (gen) {%
    \textbf{Backend emission}\\
    CUDA macros or schedule parameters\\
    memory-placement directives\\
    generated GPU kernel}; \\
};

\node[title, anchor=center]
      at (backendboxes.north |- inputtitle.center)
      {Backend compiler};

\draw[arr, rounded corners=1.5pt]
  (solve.east) -- ++(5mm,0) |- (out.west);

\draw[arr] (out.south) -- (gen.north);

\end{tikzpicture}}
\caption{Workflow of \textsc{Mosaic}. An external compiler or scheduler supplies
schedule metadata, while the target profile supplies architecture limits and
effective costs. \textsc{Mosaic} classifies dimensions, enumerates placements,
constructs candidates, filters feasibility, models cost, and performs MILP+LP
selection. The backend consumes the selected tile sizes, block dimensions, and
memory-placement directives to emit GPU code; no candidate kernel is generated
or timed during selection.}
\Description{Compact workflow diagram separating compiler inputs, the Mosaic
analytical engine, and backend code emission.}
\label{fig:mosaic-pipeline}
\end{figure*}

\name is realized as a backend-agnostic Python pipeline that consumes schedule
information, runs the analytical selection of \autoref{sec:model} outside the
compiler, and returns the chosen configuration as code-generation parameters.
It treats the schedule as an abstraction layer rather than reaching into
compiler internals: from it \name extracts the loop flags, affine access maps,
tile footprints, read/write polarities, and operation counts required by the
model, from which the program-derived quantities \(\dim(A)\),
\(K_{\mathit{total}}\), tile footprints, and stencil halo widths \(h_d\) follow
directly. The same engine drives both backends and reads identical fields from
each: PPCG exposes this information through schedule trees, and TVM through
schedule traces. For every placement scheme it assembles the Stage-1 MILP and
Stage-2 LP of \autoref{eq:two-level}, solved with PuLP and the CBC backend, and
keeps the lowest-cost scheme. The selected tiles, block dimensions, and staged
operands are emitted as backend configuration parameters, including shared memory
buffers with cooperative loads and barriers for PPCG, schedule parameters for
TVM, while capacity parameters come from device properties and the effective
\(L\)/\(c\) constants from the microbenchmarks of
\autoref{sec:gpu-target-profiling}. Because all resource reasoning happens
outside the backend, no candidate kernel is generated or timed during selection.
